\chapter{Introduction}
\section{Chapter overview}
State the flow of this chapter using the TTT approach: context, problem definition, motivation, gap, contributions, challenges, research questions, aim, and objectives.

\section{Problem background}
\subsection{Context of the research domain}

Analysing company shares has become an increasingly complex and expertise-driven task in modern financial markets. Effective evaluation requires investors to review numerous detailed documents such as annual reports, quarterly statements, regulatory announcements, and news articles, many of which contain highly technical information. Interpreting these materials demands a solid understanding of accounting and finance, skills that many individual investors lack. Even for those equipped with the necessary knowledge, analysing many interrelated parameters mentally remains a challenging and time-consuming process. 

Another major challenge arises from the contextual nature of financial markets. Analytical models developed for one country’s market often fail to produce accurate results when applied to another. Market behaviour is influenced by a range of country-specific factors, including economic conditions, availability of resources, governance structures, and political stability. Therefore, analytical models must be adapted to reflect the unique characteristics of each national market. 

Given these complexities, most high-net-worth individuals and brokerage firms employ professional analysts to conduct market evaluations manually. In contrast, retail investors often lack the time, expertise, and resources to perform such detailed analysis. As a result, they become more susceptible to misinformation, market manipulation, and financial losses. 

Retail investors frequently rely on research and recommendations provided by brokerage firms. However, over time, it has become evident that such research cannot always be considered objective or reliable. Brokerage firms often face inherent conflicts of interest, as their revenue tends to increase with higher trading volumes. Consequently, investors may be encouraged to trade more frequently, even when such activity is not in their best financial interest. 

This imbalance in knowledge and analytical capacity contributes to market inefficiencies. Due to the limited financial literacy within the investor community, certain companies may trade at unrealistically high valuations, while others remain undervalued conditions that are detrimental to both market health and broader economic stability. 

Furthermore, brokerage firms often issue target prices for companies, but forecasting such figures is inherently uncertain and difficult. Effective share market analysis should instead prioritize a deep understanding of the underlying business fundamentals rather than relying solely on externally set price targets. Before projecting future performance, it is essential to examine a company’s historical financial data and assess its current operational status. 

Although there has been substantial attention on technical analysis, the development and accessibility of robust fundamental analysis models remain limited. 

\subsection{Subdomain and key technological context}
Narrow the focus to the specific technical area that your project addresses.

\subsection{Practical and real-world context}
Explain application context, adoption constraints, and current industry or regulatory relevance.

\section{Problem definition}

\subsection{Problem statement}
Fundamental analysis is widely used in practice but remains relatively under-researched, primarily due to its reliance on large volumes of unstructured data from diverse sources. This makes it difficult for human analysts to systematically process and interpret information, highlighting the need for an Agentic AI-based approach to enable more efficient and consistent analysis. 

\section{Research motivation}

The analysis of equity investments in the Sri Lankan capital market remains a complex and often inaccessible task for many investors, particularly those with limited financial domain knowledge. At present, investors primarily rely on general-purpose large language model (LLM) tools and research reports provided by broker firms. While these resources offer some level of guidance, they are not specifically tailored to the nuances of stock market analysis, resulting in inefficiencies and a steep learning curve for users. 

Consequently, a significant number of investors are prone to misinterpretation of information and suboptimal decision-making, leading to potential financial losses. This challenge is especially pronounced among retail investors, who may lack the expertise required to critically evaluate financial data and insights. Moreover, these limitations can negatively impact overall confidence and trust in the Sri Lankan capital market. 

Therefore, there is a clear need for a more specialized, accessible, and intelligent solution that can support investors in making informed decisions, while also enhancing transparency and trustworthiness within the market ecosystem. 

\section{Research gap}
\textbf{Gap 1:} Add citation-backed explanation.

\textbf{Gap 2:} Add citation-backed explanation.

\section{Contribution to the body of knowledge}
\subsection{Contribution to the research domain}
Explain methodological or theoretical contribution.

\subsection{Contribution to the problem domain}
Explain practical or applied contribution.

\section{Research challenges}
\textbf{Challenge 1:} Describe challenge and justification.

\textbf{Challenge 2:} Describe challenge and justification.

\section{Research questions}
\resq{1}{What key parameters can be extracted using AI to improve the efficiency and effectiveness of fundamental analysis?}
\resq{2}{What is the influential impact (in detail and color context) of a multi-scale feature
extraction mechanism on enhancing low-light images in non-uniform lighting conditions?}
\resq{3}{In what ways can the GAN-based networks be directed to produce robust and generalized
enhancements across a diverse non-uniformity range of low-light scenes?}

\section{Research aim}
State the aim in one sentence using design-develop-evaluate language, then elaborate briefly.

\section{Research objectives}
\begin{itemize}
    \item RO1: Insert SMART objective 1.
    \item RO2: Insert SMART objective 2.
    \item RO3: Insert SMART objective 3.
    \item RO4: Insert SMART objective 4.
\end{itemize}

\section{Chapter summary}
Summarize key points and transition to Chapter 2.
